% !TEX root = ../../ctfp-print.tex

\lettrine[lhang=0.17]{以}{前見たように、}圏$\cat{C}$の対象$a$を固定すると、写像$\cat{C}(a, -)$は$\cat{C}$から$\Set$への（共変）関手である。
\[x \to \cat{C}(a, x)\]
（コドメインが$\Set$なのは、hom集合$\cat{C}(a, x)$が\emph{集合}だからだ。）この写像をhom関手と呼ぶ。射への作用もすでに定義した。

次に、この写像における$a$を変化させてみよう。すると、hom\emph{関手}$\cat{C}(a, -)$を任意の$a$に対応させる新しい写像が得られる。
\[a \to \cat{C}(a, -)\]
これは圏$\cat{C}$の対象から関手への写像であり、関手は関手圏（
\hyperref[natural-transformations]{自然変換}
の節を参照）の\emph{対象}である。$\cat{C}$から$\Set$への関手圏を$[\cat{C}, \Set]$と表記する。また、hom関手が典型的な
\hyperref[representable-functors]{表現可能関手}
であることも思い出してほしい。

二つの圏の間の対象の写像があるとき、その写像が関手でもあるかどうかを問うのは自然なことだ。つまり、一方の圏の射を他方の圏の射に持ち上げられるかどうかということだ。$\cat{C}$の射は単に$\cat{C}(a, b)$の要素だが、関手圏$[\cat{C}, \Set]$の射は自然変換である。したがって、射を自然変換に対応させる写像を探すことになる。

射$f \Colon a \to b$に対応する自然変換を見つけられるか考えてみよう。まず、$a$と$b$がそれぞれ何に写されるかを確認する。それらは二つの関手$\cat{C}(a, -)$と$\cat{C}(b, -)$に写される。これら二つの関手の間の自然変換が必要だ。

ここで米田の補題を使うのがポイントだ：
\[[\cat{C}, \Set](\cat{C}(a, -), F) \cong F a\]
一般的な$F$をhom関手$\cat{C}(b, -)$で置き換えると：
\[[\cat{C}, \Set](\cat{C}(a, -), \cat{C}(b, -)) \cong \cat{C}(b, a)\]
\begin{figure}[H]
  \centering
  \includegraphics[width=0.6\textwidth]{images/yoneda-embedding.jpg}
\end{figure}

\noindent
これはまさに探し求めていた二つのhom関手の間の自然変換だが、少しひねりがある。自然変換と射の対応において、射は$\cat{C}(b, a)$の要素、つまり「逆方向」に向かうのだ。しかしこれは問題ない。見ている関手が反変であることを意味するだけだ。

\begin{figure}[H]
  \centering
  \includegraphics[width=0.65\textwidth]{images/yoneda-embedding-2.jpg}
\end{figure}

\noindent
実は、期待以上のものが得られた。$\cat{C}$から$[\cat{C}, \Set]$への写像は反変関手であるだけでなく、\emph{完全忠実}関手でもある。完全性と忠実性は、関手がhom集合をどのように写すかを表す性質だ。

\emph{忠実}関手はhom集合上で\newterm{単射}（injective）である。すなわち、異なる射を異なる射に写す。射を合併することはない。

\emph{充満}関手はhom集合上で\newterm{全射}（surjective）である。すなわち、あるhom集合を別のhom集合の\emph{上に}写し、後者を完全に被覆する。

完全忠実関手$F$はhom集合上で\newterm{全単射}（bijection）である。つまり、両方の集合のすべての要素の一対一対応だ。元の圏$\cat{C}$の任意の対象の対$a$と$b$に対して、$\cat{C}(a, b)$と$\cat{D}(F a, F b)$の間に全単射が存在する。ここで$\cat{D}$は$F$の行先の圏である（今回の場合は関手圏$[\cat{C}, \Set]$）。これは$F$が\emph{対象}上で全単射であることを意味しないことに注意しよう。$\cat{D}$には$F$の像に含まれない対象が存在する可能性があり、そのような対象のhom集合については何も言えない。

\section{埋め込み}

今説明した（反変）関手、つまり$\cat{C}$の対象を$[\cat{C}, \Set]$の関手に写す関手：
\[a \to \cat{C}(a, -)\]
は\newterm{米田埋め込み}（Yoneda embedding）を定義する。これは圏$\cat{C}$（厳密には反変性のため圏$\cat{C}^\mathit{op}$）を関手圏$[\cat{C}, \Set]$の中に\emph{埋め込む}。$\cat{C}$の対象を関手に写すだけでなく、それらの間のすべての接続を忠実に保持する。

これは非常に有用な結果だ。数学者たちは関手の圏、特にコドメインが$\Set$の関手の圏について多くの知識を持っているからだ。任意の圏$\cat{C}$を関手圏に埋め込むことで、多くの洞察を得られる。

もちろん、米田埋め込みの双対版もある。\newterm{余米田埋め込み}（co-Yoneda embedding）と呼ばれることもある。各hom集合のソース対象ではなくターゲット対象を固定することから始めることができる。$\cat{C}(-, a)$だ。これにより反変hom関手が得られる。$\cat{C}$から$\Set$への反変関手は馴染み深い前層（\hyperref[limits-and-colimits]{極限と余極限}を参照）だ。余米田埋め込みは圏$\cat{C}$の前層の圏への埋め込みを定義する。射への作用は次の式で与えられる：
\[[\cat{C}^\mathit{op}, \Set](\cat{C}(-, a), \cat{C}(-, b)) \cong \cat{C}(a, b)\]
ここでも数学者たちは前層の圏について多くの知識を持っているため、任意の圏をその中に埋め込めることは大きな利点だ。

\section{Haskellへの応用}

Haskellでは、米田埋め込みはreaderファンクタ間の自然変換と逆方向の関数の間の同型として表現できる：

\begin{snipv}
forall x. (a -> x) -> (b -> x) \ensuremath{\cong} b -> a
\end{snipv}
（readerファンクタは\code{((->) a)}と同値であることを思い出そう。）

この等式の左辺は多相関数で、\code{a}から\code{x}への関数と型\code{b}の値が与えられたとき、型\code{x}の値を生成できる（関数\code{b -> x}の括弧を外してカリー化を解除している）。これをすべての\code{x}に対して行えるのは、その関数が\code{b}を\code{a}に変換する方法を知っている場合だけだ。秘密裏に関数\code{b -> a}にアクセスできなければならない。

そのような変換関数\code{btoa}が与えられると、左辺の関数（\code{fromY}と呼ぼう）を次のように定義できる：

\src{snippet01}
逆に、関数\code{fromY}が与えられると、恒等関数を渡して呼び出すことで変換関数を復元できる：

\src{snippet02}
これにより、\code{fromY}型の関数と\code{btoa}の間の全単射が確立される。

この同型の別の見方として、これは\code{b}から\code{a}への関数の\acronym{CPS}変換（continuation-passing style）だということがある。引数\code{a -> x}は継続（ハンドラ）である。結果は\code{b}から\code{x}への関数で、型\code{b}の値で呼び出されたとき、エンコードされる関数が前合成された継続を実行する。

米田埋め込みはHaskellのデータ構造の代替表現の一部も説明する。特に、\code{Control.Lens}ライブラリのレンズの
\urlref{https://bartoszmilewski.com/2015/07/13/from-lenses-to-yoneda-embedding/}
{非常に有用な表現}
を提供する。

\section{前順序の例}

この例はRobert Harperによって提案された。前順序によって定義される圏への米田埋め込みの応用だ。\newterm{前順序}（preorder）は、要素間に順序関係を持つ集合で、その関係は伝統的に$\leqslant$（以下）と書かれる。「前」という接頭語がついているのは、この関係が推移的かつ反射的であることだけを要求しており、反対称性は必ずしも要求されないためだ（サイクルが存在する可能性がある）。

前順序関係を持つ集合は圏を生み出す。対象はこの集合の要素だ。対象$a$から$b$への射は、対象が比較できない場合や$a \leqslant b$でない場合は存在せず、$a \leqslant b$の場合は$a$から$b$に向かう射が存在する。一つの対象から別の対象への射は高々一つしかない。したがって、そのような圏のhom集合は空集合か一要素集合のどちらかだ。そのような圏を\emph{薄い}（thin）圏と呼ぶ。

この構成が確かに圏であることは容易に確認できる。$a \leqslant b$かつ$b \leqslant c$ならば$a \leqslant c$なので矢印は合成可能であり、合成は結合的だ。また、すべての要素は自分自身と等しい（以下である）ので（基礎関係の反射性）、恒等射も存在する。

次に、この前順序圏に余米田埋め込みを適用しよう。特に射への作用に注目する：
\[[\cat{C}, \Set](\cat{C}(-, a), \cat{C}(-, b)) \cong \cat{C}(a, b)\]
右辺のhom集合は$a \leqslant b$の場合に限り空でない（その場合は一要素集合だ）。
したがって、$a \leqslant b$のとき、左辺に唯一の自然変換が存在する。そうでなければ自然変換は存在しない。

では、前順序におけるhom関手間の自然変換とは何だろうか？それは集合$\cat{C}(-, a)$と$\cat{C}(-, b)$の間の関数の族であるべきだ。前順序では、これらの各集合は空か単集合のどちらかだ。どのような関数が利用可能か見てみよう。

空集合から空集合への関数（空集合上の恒等）、空集合から単集合への関数\code{absurd}（空集合の要素について定義するだけでよいので何もしない）、単集合から単集合への関数（一要素集合上の恒等）がある。唯一禁じられているのは単集合から空集合への写像だ（そのような関数が唯一の要素に作用したときの値は何になるのだろうか？）。

したがって、この自然変換は単集合のhom集合を空のhom集合に接続することはない。言い換えると、$x \leqslant a$（単集合のhom集合$\cat{C}(x, a)$）ならば$\cat{C}(x, b)$は空になれない。$\cat{C}(x, b)$が空でないということは$x$が$b$以下であることを意味する。したがって、問題の自然変換の存在は、すべての$x$について$x \leqslant a$ならば$x \leqslant b$であることを要求する。
\[\text{すべての } x \text{ について}, x \leqslant a \Rightarrow x \leqslant b\]
一方、余米田は、この自然変換の存在が$\cat{C}(a, b)$が空でないこと、つまり$a \leqslant b$と同値であることを教えてくれる。合わせると：
\[a \leqslant b \text{ であることと、すべての } x \text{ について } x \leqslant a \Rightarrow x \leqslant b \text{ であることは同値}\]
この結果には直接到達することもできた。直感的には、$a \leqslant b$ならば$a$より下にある要素はすべて$b$より下にもなければならない。逆に、右辺の$x$に$a$を代入すると$a \leqslant b$が導かれる。しかし米田埋め込みを通じてこの結果に到達する方がはるかに刺激的だと認めなければならない。

\section{自然性}

米田の補題は自然変換の集合と$\Set$の対象の間の同型を確立する。自然変換は関手圏$[\cat{C}, \Set]$の射だ。任意の二つの関手の間の自然変換の集合はその圏のhom集合だ。米田の補題の同型は：
\[[\cat{C}, \Set](\cat{C}(a, -), F) \cong F a\]
この同型は$F$と$a$の両方について自然であることがわかっている。言い換えると、積圏$[\cat{C}, \Set] \times \cat{C}$から取られた対$(F, a)$について自然だ。ここで$F$は関手圏の\emph{対象}として扱っていることに注意しよう。

これが何を意味するか少し考えてみよう。自然同型とは二つの関手の間の可逆な\emph{自然変換}だ。そして確かに、同型の右辺は関手だ。それは$[\cat{C}, \Set] \times \cat{C}$から$\Set$への関手であり、対$(F, a)$への作用は集合、つまり対象$a$で関手$F$を評価した結果だ。これを\newterm{評価関手}（evaluation functor）と呼ぶ。

左辺もまた$(F, a)$を自然変換の集合$[\cat{C}, \Set](\cat{C}(a, -), F)$に写す関手だ。

これらが本当に関手であることを示すには、射への作用も定義する必要がある。では、対$(F, a)$と$(G, b)$の間の射は何だろうか？それは射の対$(\Phi, f)$であり、最初のものは関手間の射（自然変換）で、二番目は$\cat{C}$の通常の射だ。

評価関手はこの対$(\Phi, f)$を取り、二つの集合$F a$と$G b$の間の関数に写す。$a$における$\Phi$の成分（$F a$を$G a$に写す）と$G$によって持ち上げられた射$f$から、そのような関数を簡単に構成できる：
\[(G f) \circ \Phi_a\]
$\Phi$の自然性により、これは次と同じだ：
\[\Phi_b \circ (F f)\]
同型全体の自然性を証明するつもりはない。関手が何であるかを確立した後は、証明はかなり機械的だ。同型が関手と自然変換から構成されているという事実から従う。それが崩れる方法はないのだ。

\section{練習問題}

\begin{enumerate}
  \tightlist
  \item
        余米田埋め込みをHaskellで表現せよ。
  \item
        \code{fromY}と\code{btoa}の間で確立した全単射が同型（二つの写像が互いの逆である）であることを示せ。
  \item
        モノイドに対する米田埋め込みを導出せよ。モノイドの単一の対象にはどの関手が対応するか？モノイド射にはどの自然変換が対応するか？
  \item
        前順序への\emph{共変}米田埋め込みの応用は何か？（Gershom Bazermanによる問い）
  \item
        米田埋め込みを使って任意の関手圏$[\cat{C}, \cat{D}]$を関手圏$[[\cat{C}, \cat{D}], \Set]$に埋め込むことができる。これが射（この場合は自然変換）に対してどのように作用するかを明らかにせよ。
\end{enumerate}
